\chapter{Исследовательская часть}

\section{Параметризация муравьиного алгоритма}

% по итогам выполненного исследования дать рекомендации о значениях либо диапазонах значений параметров, настройках и/или режимах работы метода, основанного на муравьином алгоритме;

Результаты параметризации муравьиного алгоритма по параметрам $\rho$, $\alpha$ и $t_{max}$ приведены в таблице \ref{par_table}.
По полученным данным $t_{max}$ должен лежать в диапазоне [200, 400], а $\rho$ -- [0.1, 0.25]. Значение $\alpha$ не оказывает особое влияние на итоговый результат.
Наиболее оптимальным является следующий набор:
\begin{itemize}
	\item $\rho = 0.25$;
	\item $\alpha = 0.9$;
	\item $t_{max} = 200$.
\end{itemize}

% тут должны быть графики по таблице

\section{Замеры реального времени}

Замеры проводились на ноутбуке Redmibook Pro 14 2022:
\begin{itemize}
	\item процессор Intel i7-12650H (архитектура x86-64, 10 ядер по 2.3 ГГц);
	\item оперативная память 16 Гб (DDR4, 5200 МГц);
	\item операционная система Windows 11 Home 25H2.
\end{itemize}
Для замеров использовалась функция now() из библиотеки chrono~\cite[с.~705]{lit1}, результаты представлены на рис.~4.1.

\begin{figure}[!htbp]
	\centering
	\begin{tikzpicture}
		\selectcolormodel{gray}
		\begin{semilogyaxis}[
			axis lines = left, xlabel = Размер, ylabel = {Время, с}, grid=both,
			minor grid style={gray!25}, major grid style={gray!50}, 
			legend style={font=\scriptsize, legend pos=north west}, 
			ymin=5e-7, ymax=5e4, xmin=1.5, xmax=15.5
		]
			\addplot+[smooth, mark=none] table [x index=0, y index=1] {../code/timeres.txt};
			\addlegendentry{Алгоритм полного перебора}
			\addplot+[smooth, mark=none, dashed] table [x index=0, y index=2] {../code/timeres.txt};
			\addlegendentry{Муравьиный алгоритм}
		\end{semilogyaxis}
	\end{tikzpicture}
	\caption{Графики зависимости времени выполнения от размера графа.}
\end{figure}

\section*{Вывод}

В данном разделе была проведена параметризация муравьиного алгоритма и получена наиболее оптимальная комбинация его параметров. Также были проведены замеры времени работы для каждой реализации из предыдущего раздела. Полученные на основе этих замеров графики подтверждают рассчитанные ранее трудоёмкости рассматриваемых в работе алгоритмов.

\clearpage
